% ECONOMICS_PROBLEM: {"id": "I11-628", "title": "Realized returns to medical AI", "jel_code": "I11", "source_id": "OP-0447"}
% !TeX program = xelatex
\documentclass[11pt,a4paper]{article}
\usepackage[margin=23mm]{geometry}
\usepackage{fontspec}
\setmainfont{Latin Modern Roman}
\usepackage{amsmath,amssymb,mathtools}
\usepackage{xcolor,enumitem,booktabs,longtable,array,xurl,hyperref}
\definecolor{muted}{HTML}{53636C}
\hypersetup{colorlinks=true,urlcolor=blue,linkcolor=blue}
\setlength{\parindent}{0pt}
\setlength{\parskip}{5pt}
\newcommand{\R}{\mathbb R}
\newcommand{\E}{\mathbb E}
\newcommand{\N}{\mathbb N}
\newcommand{\ind}{\mathbf 1}
\newcommand{\Var}{\operatorname{Var}}
\newcommand{\Cov}{\operatorname{Cov}}
\newcommand{\supp}{\operatorname{supp}}
\DeclareMathOperator*{\argmin}{arg\,min}
\DeclareMathOperator*{\argmax}{arg\,max}
\newcommand{\entrymeta}[3]{{\sffamily\small\color{muted}\textbf{\texttt{#1}}\quad|\quad #2\par #3\par}}
\newcommand{\Problemref}[1]{\texttt{#1}}
\newcommand{\JELref}[2]{\texttt{#2} (JEL \texttt{#1})}
\begin{document}
\section*{Realized returns to medical AI}
\textbf{Problem ID:} \texttt{I11-628}\quad \textbf{JEL:} \texttt{I11}\par
% BEGIN ECONOMICS STATEMENT
\label{problem:OP-0447}
{\sffamily\small\color{muted}Original subject group: Health economics\par}
\label{definitions:problem:30007729}
\textbf{Audit classification:} Background; proposed extension. Prevented training/supervision cost double counting and post-treatment visit selection.
\subsection*{Definitions and precise research target}
\textit{Editorial formalization.} Consider adult primary-care visits in a prespecified national provider network over two years. Randomize clinics to access \(a=1\) to one version-controlled artificial-intelligence system for documentation and decision support, or \(a=0\) to existing practice. Let \(Q_i(a)\) be quality-adjusted health, \(C_i(a)\) total real care costs, and \(L_i(a)\) clinician labor hours associated with patient \(i\). For health valuation \(\lambda>0\), define
\[\tau_W=E[\lambda\{Q_i(1)-Q_i(0)\}-\{C_i(1)-C_i(0)\}]-k,\qquad\tau_L=E[L_i(1)-L_i(0)],\]
where \(k\) is implementation, maintenance, supervision, and training resource cost per eligible patient. Costs include follow-up testing, downstream hospital use, and human review of errors. Determine realized welfare and labor effects under randomized clinic assignment, not savings implied by automatable task shares. Track uptake, diagnostic mistakes, altered visit volume, and effects by patient complexity. Preserve baseline patient inclusion and specify clinic spillovers and model-version changes. Longer-term organizational responses require a separate extension or explicit extrapolation assumptions.

The unit of assignment is the clinic; the treatment is an offered deployment package including software version, allowed tasks, training, review, and fallback rules. Count \(Q_i\) in quality-adjusted life-years over two years and \(L_i\) in all associated clinician hours, including error correction and downstream visits. Allocate fixed deployment resources across a declared denominator of baseline eligible patients to obtain \(k\). If training or supervision labor enters \(k\), exclude the identical cost from \(C_i\); otherwise include all labor opportunity costs in \(C_i\) and let \(k\) contain only remaining resources. Randomized offers identify adoption-inclusive mean contrasts under consistency and no interference. If visit volume changes, a baseline patient-cohort outcome or per-clinic total is appropriate; a post-treatment visit average compares selected populations. All expectations use a stated baseline population and probability law. Identification means every admissible data-generating model with the same observed-data law yields the same displayed target; otherwise report the identified set. Exact equations, horizons and assumptions are editorial, rather than source-given canonical conjectures.
\subsection*{Variants and assumption changes}
\begin{enumerate}
\item \textbf{Editorial organizational variant.} Extend follow-up to ten years allowing staffing and visit volume to adjust. Report baseline-cohort and clinic totals separately.
\item \textbf{Editorial updating variant.} Compare a frozen system with a prespecified update rule based only on past data. Updating is part of the assigned intervention.
\end{enumerate}
\subsection*{References and source audit}
\textbf{1.} Projected spending savings provide context, not the trial estimand or proof of absence of causal studies. Locator: NBER Working Paper 30857; abstract. Primary indexed record and author publication page checked; direct NBER retrieval failed. URL: \url{https://www.nber.org/papers/w30857}.
\begin{enumerate}\raggedright
\item\hypertarget{definitions:source:30007729:1}{} Nikhil Sahni; George Stein; Rodney Zemmel; David M. Cutler. 2023. \emph{The Potential Impact of Artificial Intelligence on Healthcare Spending}. NBER Working Paper 30857; abstract.
\url{https://www.nber.org/papers/w30857}.
DOI: \url{https://doi.org/10.3386/w30857}.
\end{enumerate}

\subsection*{Common conventions: Source mathematical conventions}

These conventions apply to each entry unless explicitly replaced.

\textbf{Models and identification.} A primitive model \(m\) specifies all probability laws, feasible choices, information, equilibrium and measurement rules used by its target. The admissible class \(\mathcal M\) is fixed independently of outcomes. If \(P_O(m)\) is its observable probability law and \(\tau(m)\) the target, its identified set at observable law \(P\) is
\[
 \mathcal I_\tau(P)=\{\tau(m):m\in\mathcal M,\ P_O(m)=P\}.
\]
Point identification means that this set is a singleton. An empty set signals incompatibility of data and assumptions. Coordinate bounds need not describe the joint set. Bounds are sharp only if every retained target is attainable or arbitrarily approachable by compatible models. Finite bounds require explicit restrictions on unobserved quantities; unrestricted confounding, hidden wealth, extrapolation or arbitrary equilibrium selection can yield uninformative sets.

\textbf{Causal interventions.} A potential outcome \(Y_i(p)\) is the outcome under a complete policy regime \(p\), including timing, financing, information and market spillovers. The target population and units are fixed before treatment. With interference, use the complete assignment vector or a justified exposure mapping; individual no-interference notation alone cannot identify an economy-wide intervention. A difference of mean potential outcomes does not require the cross-world joint distribution. Conditioning on post-treatment migration, survival, enrollment or employment changes the estimand and needs separate justification.

\textbf{Statistical statements.} An estimator, confidence set, forecast or test specifies a sampling law, model class and loss/coverage criterion. Uniform \((1-\alpha)\)-coverage means \(\inf_{m\in\mathcal M}\Pr_m\{\tau(m)\in C_n\}\geq1-\alpha\), or an explicitly asymptotic version. The whole parameter space trivially has coverage, so informativeness requires a stated width, risk or power objective. A forecasting improvement is relative to a named benchmark, information set and evaluation population.

\textbf{Equilibrium and optimization.} An equilibrium comprises feasible plans, optimality, beliefs where needed, budget constraints and all market/resource clearing. The primitives or a stated benchmark define each correspondence \(\mathcal E(p;m)\). Nonemptiness is assumed or proved, never inferred just from compact policy sets. A selected-equilibrium target uses a declared selection rule; a robust target quantifies over all permitted equilibria. Use suprema or infima unless attainment is established. An \(\varepsilon\)-optimal policy is feasible and within \(\varepsilon\) of the finite optimum value. Report an empty feasible set separately.

\textbf{Welfare and accounting.} Normative utility, interpersonal normalization, social weights, numeraire, discounting and the use of public receipts are primitives. Behavioral utility need not coincide with the welfare criterion. Household welfare includes ownership income and rents once; adding profits again to compensating variation generally double-counts them. A monetary equivalent is defined by an explicit transfer and price convention, with monotonicity and range conditions. Average effects, mechanism contributions, transfers and welfare losses are different targets. Infinite sums require convergence and dynamic feasible plans require terminal or no-Ponzi conditions.

\textbf{Source versions.} A working-paper date, revision date and journal publication year are distinguished. A DOI for a journal version is not automatically the DOI of an earlier manuscript. Locators refer to the stated version and printed pages where specified. A metadata-only check does not verify a claim about a full-text conclusion. Every entry's source subsection narrows the provenance claim accordingly.
% END ECONOMICS STATEMENT
\end{document}
